This section studies learning problems on the Boolean cube \(\mathbb{B}^n\) through the lens of symmetry, closure, and compositional structure. The guiding question is not merely how to fit a function \(f : \mathbb{B}^n \to \mathbb{R}\) or \(f : \mathbb{B}^n \to \mathbb{B}^m\), but how to do so in a way that respects the invariances of the task class and the algebraic constraints inherited from the surrounding system. In the present book, this means treating symmetry-aware learning as another instance of closure under admissible transformations: the hypothesis class should be stable under the relevant group action, and the learned representation should expose the same structural regularities that appear in circuits, clones, and compositional hierarchies.

A useful starting point is the generalization error under a symmetry class, denoted here by \(\varepsilon(r)\), where \(r\) indexes the amount of symmetry or the radius of approximation under the chosen metric. Informally, \(\varepsilon(r)\) measures how much predictive error remains when the learner is constrained to a family of models that is invariant, or approximately invariant, under a prescribed symmetry group. In the idealized case, the target function lies in a symmetry class \(\mathcal{C}\) and the hypothesis class \(\mathcal{H}\) is closed under the same action; then the effective sample complexity can improve because the learner need not rediscover equivalent configurations separately. In practice, however, symmetry constraints can also introduce bias if the true target only partially respects the assumed invariance. The role of \(\varepsilon(r)\) is therefore to make explicit the tradeoff between structural restriction and approximation power.

Fourier analysis on \(\mathbb{B}^n\) provides a canonical language for this tradeoff. Writing Boolean functions in the Walsh--Fourier basis decomposes them into coefficients indexed by subsets of coordinates, thereby separating low-order interactions from higher-order ones. This decomposition is especially useful when the target exhibits junta or sparsity structure. A junta depends on only a small number of coordinates, while a sparse Fourier spectrum concentrates most of its mass on a small number of coefficients. Both conditions are compatible with symmetry-aware learning: if a function is invariant under a permutation group, then its Fourier support is organized into orbits, and the effective number of distinct parameters may be much smaller than the ambient dimension suggests. Consequently, learning algorithms can exploit coordinate sparsity, spectral sparsity, or orbit sparsity, depending on which structural assumption is most faithful to the data.

The symmetry groups of interest are not limited to simple coordinate permutations. Group-equivariant networks generalize the familiar idea of convolutional weight sharing to arbitrary group actions, including permutation groups and wreath products. A network is equivariant when applying a group element to the input induces a corresponding, predictable transformation of the output. This property is valuable because it enforces consistency across symmetric instances and reduces the burden on the optimizer. For example, if a task is invariant under relabeling of inputs, then a permutation-equivariant architecture can represent the same decision rule with fewer free parameters and with better inductive bias than an unconstrained multilayer perceptron. Wreath products arise naturally when the symmetry has a hierarchical or block-structured form, such as permutations within local neighborhoods together with permutations of the neighborhoods themselves. In such cases, the architecture should reflect the semidirect layering of the symmetry rather than flattening it into a single undifferentiated action.

From the perspective of closure, the central design question is whether the chosen constraints remain stable across layers. If each layer is equivariant with respect to a group action, then the composition of layers is equivariant as well, provided the intermediate nonlinearities and pooling operations are compatible with the action. This closure across layers is not automatic: a pointwise nonlinearity may preserve equivariance for some actions but not for others, and a normalization step may inadvertently break identifiability by collapsing distinct symmetric states. The architecture must therefore be checked as a whole, not merely layer by layer in isolation. In a compositional setting, one asks whether the class of admissible layers forms a closed family under the relevant operations, and whether the resulting model class remains identifiable up to the intended symmetry. Identifiability here means that equivalent parameterizations should correspond exactly to the same function modulo the symmetry, rather than introducing spurious degrees of freedom that obscure interpretation or hinder optimization.

This closure perspective also clarifies the relationship between representation learning and circuit extraction. If a learned model is constrained to respect a symmetry class, then its internal structure may be distilled into a smaller circuit or symbolic artifact that preserves the same invariance profile. Such distillation is especially attractive when the learned model is used as an intermediate artifact in a larger verified pipeline. The resulting model zoo can be organized by symmetry type, approximation regime, and compression level, allowing one to compare architectures not only by accuracy but also by equivariance, parameter efficiency, and extractability. In this setting, the artifact is not merely a trained network but a family of related representations: a dense model, a symmetry-restricted model, and a distilled circuit approximation that can be checked against the original behavior on a chosen test suite.

Privacy considerations enter naturally once symmetry-aware learning is deployed on sensitive data. Symmetry can reduce the effective hypothesis space, which may improve generalization, but it can also make certain attributes easier to infer if the symmetry class aligns with latent structure in the data. Moreover, distillation and model compression can leak information if the extracted artifact preserves too much of the training signal. A careful treatment therefore requires distinguishing between invariance that is beneficial for robustness and invariance that inadvertently exposes protected correlations. In practical terms, one should evaluate whether the symmetry constraints amplify membership inference, attribute inference, or reconstruction risk, especially when the model zoo includes distilled circuits or other compact artifacts intended for downstream reuse. Privacy-aware design may require additional regularization, noise injection, or access controls on the exported representations.

The overall lesson is that learning on \(\mathbb{B}^n\) is most effective when the hypothesis class is aligned with the algebraic and symmetry structure of the task. Fourier methods reveal which interactions matter; junta and sparsity assumptions identify low-complexity targets; equivariant architectures encode invariance directly; and closure across layers ensures that the intended symmetry survives composition. At the same time, identifiability and privacy must be treated as first-class constraints, not afterthoughts. In the language of this book, a successful learning architecture is one whose symmetry constraints are compositional, whose representations are interpretable up to the intended equivalence, and whose exported artifacts remain faithful to both performance and governance requirements.
